% !TeX root = tcolorbox.tex
% include file of tcolorbox.tex (manual of the LaTeX package tcolorbox)
\clearpage
\section{逐字文本的保存与读取}\label{sec:verbatim}%
\tcbset{external/prefix=external/verbatim_}%
以下宏是对已知宏包 \refPkg{moreverb} 与 \refPkg{verbatim} 中原有宏的小幅修改版本。
它们会在 |tcolorbox| 环境内部被隐式使用，但也可在外部使用。
\enlargethispage*{1.5cm}

\begin{docEnvironment}{tcbverbatimwrite}{\marg{file name}}
  将 \meta{环境内容}保存到由 \meta{文件名}指定的文件中。
  环境内部的 \TeX\ 宏不会被展开。
\begin{dispExample}
\begin{tcbverbatimwrite}{\jobname_verbexp.tex}
  This text is saved \textit{as is}.
\end{tcbverbatimwrite}

Now, we are using the file:\par
\input{\jobname_verbexp.tex}
\end{dispExample}

本环境可在自定义环境内部使用。注意，在环境定义内部，
必须使用 |\tcbverbatimwrite| 而非 |\begin{tcbverbatimwrite}|，
并使用 |\endtcbverbatimwrite| 而非 |\end{tcbverbatimwrite}|。
\begin{dispExample}
\newenvironment{myverbatim}{%
  \begingroup\tcbverbatimwrite{\jobname_myverb.tex}}%
  {\endtcbverbatimwrite\endgroup}

\begin{myverbatim}
  This is the text which is saved by my own environment.
\end{myverbatim}

Now, we are using the file:\par
\input{\jobname_myverb.tex}
\end{dispExample}
\end{docEnvironment}

\begin{docEnvironment}{tcbwritetemp}{}
  功能与 \refEnv{tcbverbatimwrite} 相同，但使用 \refKey{/tcb/tempfile}
  键的值作为文件名。
\begin{dispExample}
\begin{tcbwritetemp}
  This text is saved \textit{as is}.
\end{tcbwritetemp}

Now, we are using the file:\par
\tcbusetemp
\end{dispExample}
\end{docEnvironment}


\begin{docCommand}{tcbusetemp}{}
  读取由 \refEnv{tcbwritetemp} 保存的当前临时文件。
\end{docCommand}

\clearpage

\begin{docTcbKey}[][doc new=2017-07-05]{verbatim ignore percent}{\colOpt{=true\textbar false}}{默认 |true|，初始为 |false|}
  若将本选项设为 |true|，则对于 \refEnv{tcbverbatimwrite} 以及所有基于
  \refEnv{tcbverbatimwrite} 构建的宏和环境（例如 \refEnv{tcbwritetemp}、
  \refEnv{tcblisting} 或 \refEnv{dispExample}），百分号 \% 会被静默忽略。\par
  本选项可用于创造一些特殊效果，但主要面向使用 DocStrip 编写文档的场合。
  本选项的诞生受到了 Yudai Nakata 的启发。
  注意，本选项不会被 \refKey{/tcb/reset} 重置。
\begin{dispExample*}{sidebyside}
Normal usage:
\begin{tcbwritetemp}
%\begin{center}\bfseries
This is my text.
%\end{center}
\end{tcbwritetemp}

\tcbusetemp

\tcbset{verbatim ignore percent}
\bigskip Option applied:
\begin{tcbwritetemp}
%\begin{center}\bfseries
This is my text.
%\end{center}
\end{tcbwritetemp}

\tcbusetemp
\end{dispExample*}

\begin{marker}
注意，\emph{每个}百分号都会被移除，包括已转义的百分号。
\end{marker}

\begin{dispExample*}{sidebyside}
% \tcbuselibrary{listings}
\def\percent{\%}

\begin{tcblisting}{title=Normal}
%\begin{center}\bfseries
This is my 5\percent\ text
and this is my 10\% text.
%\end{center}
\end{tcblisting}

\begin{tcblisting}{
  title=Option applied,
  verbatim ignore percent}
%\begin{center}\bfseries
This is my 5\percent\ text
and this is my 10\% text.
%\end{center}
\end{tcblisting}
\end{dispExample*}

\end{docTcbKey}

\clearpage

\begin{docTcbKey}[][doc new=2023-02-10]{verbatim ignore indention at end}{\colOpt{=true\textbar false}}{默认 |true|，初始为 |true|}
  若将本选项设为 |true|，则同一行上位于 |\end{tcbverbatimwrite}| 之前的
  空格与制表符会被忽略。所有基于 \refEnv{tcbverbatimwrite} 的环境同理。
  若同一行上 |\end{tcbverbatimwrite}| 之前还有其他内容，则会写出全部内容。\par
  若将本选项设为 |false|，则会写出 |\end{tcbverbatimwrite}| 之前的一切内容。
  这是 6.0.0 版本之前的默认行为。\par
  \refKey{/tcb/verbatim ignore indention at end} 的设计初衷是让环境可以缩进。
  本选项的诞生受到了 Yukai Chou 的启发。

\begin{dispExample}
% \tcbuselibrary{minted}
\texttt{verbatim ignore indention at end=false}

    \begin{tcblisting}{listing engine=minted, listing only,
                       verbatim ignore indention at end=false}
      int a;
    \end{tcblisting}

    \begin{tcblisting}{listing engine=minted, listing only,
                       verbatim ignore indention at end=false}
      int a;
      int b;  \end{tcblisting}

\texttt{verbatim ignore indention at end=true}

    \begin{tcblisting}{listing engine=minted, listing only,
                       verbatim ignore indention at end}
      int a;
    \end{tcblisting}

    \begin{tcblisting}{listing engine=minted, listing only,
                       verbatim ignore indention at end}
      int a;
      int b;  \end{tcblisting}
\end{dispExample}
\end{docTcbKey}
